\subsection*{Решение задачи Коши}
\addcontentsline{toc}{subsection}{Решение задачи Коши}

Дано дифференциальное уравнение второго порядка:
\[
    y'' + py' + qy = f(x), \quad p = -1, \quad q = 0, \quad f(x) = 3
\]

с начальными условиями:
\[
    y(0) = 6, \quad y'(0) = 2
\]

Исходное дифференциальное уравнение принимает вид:
\[
    y'' - y' = 3
\]

\textbf{Решение однородного уравнения:}

Характеристическое уравнение: $r(r-1) = 0$, откуда $r_1 = 0$, $r_2 = 1$.

Общее решение однородного уравнения:
\[
    y_h(x) = C_1 e^{0x} + C_2 e^{1x} = C_1 + C_2 e^x
\]

\textbf{Частное решение неоднородного уравнения:}

Так как правая часть $f(x) = 3$ \textemdash\thinspace константа, ищем частное решение в виде $y_p = Ax$.


Подставляя в исходное уравнение:
\begin{align*}
    y_p'= A, \quad y_p'' = 0 \\
    0 - A = 3 \quad \Rightarrow \quad A = -3
\end{align*}
Таким образом, $y_p(x) = -3x$.

\textbf{Общее решение:}

\[
    y(x) = C_1 + C_2 e^x - 3x
\]

Производная:
\[
    y'(x) = C_2 e^x - 3
\]

\textbf{Определение констант интегрирования:}

Из начального условия $y'(0) = 2$:
\[
    C_2 \cdot 1 - 3 = 2 \quad \Rightarrow \quad C_2 = 5
\]

Из начального условия $y(0) = 6$:
\[
    C_1 + 5 - 0 = 6 \quad \Rightarrow \quad C_1 = 1
\]

\textbf{Точное аналитическое решение задачи Коши:}
\[
    y(x) = 1 + 5e^x - 3x
\]

Вычисление граничного условия $y(1)$:
\[
    b = y(1) = 1 + 5e - 3 = 5e - 2
\]

\subsection*{Метод прогонки}
\addcontentsline{toc}{subsection}{Метод прогонки}

Рассмотрим трёхдиагональную систему:
\[
    a_i y_{i-1} + b_i y_i + c_i y_{i+1} = d_i, \quad i = 1, \ldots, m
\]

с граничными условиями $y_0 = a$ и $y_{m+1} = b$.

\textbf{Построение разностной схемы при $n = 10$:}

Для равномерной сетки на $[0,1]$ с $n = 10$ узлами:
\[
    h = \frac{1-0}{10} = 0.1, \quad x_i = ih, \quad i = 0, \ldots, 10
\]

Разностная схема для дифференциального уравнения второго порядка:
\[
    y'' - y' = 3
\]

имеет вид:
\[
    a_i y_{i-1} + b_i y_i + c_i y_{i+1} = d_i
\]

где коэффициенты при $p_i = -1$, $q_i = 0$, $f_i = 3$:
\begin{align*}
    a_i & = 1 - \frac{h}{2}p_i = 1 - \frac{0.1}{2}(-1) = 1.05 \\
    b_i & = h^2 q_i - 2 = 0 - 2 = -2                          \\
    c_i & = 1 + \frac{h}{2}p_i = 1 + \frac{0.1}{2}(-1) = 0.95 \\
    d_i & = h^2 f_i = 0.01 \cdot 3 = 0.03
\end{align*}

Для $i = 1, \ldots, 9$ система принимает вид:
\[
    1.05 y_{i-1} - 2 y_i + 0.95 y_{i+1} = 0.03
\]

с граничными условиями:
\[
    y_0 = a = 6, \quad y_{10} = b = 5e - 2
\]

\textbf{Трёхдиагональная СЛАУ для неизвестных $y_1, \ldots, y_9$:}

Для $i = 1$ подставляем $y_0 = 6$:
\[
    -2 y_1 + 0.95 y_2 = 0.03 - 6.3 = -6.27
\]

Для $i = 2, \ldots, 8$:
\[
    1.05 y_{i-1} - 2 y_i + 0.95 y_{i+1} = 0.03
\]

Для $i = 9$ подставляем $y_{10} = b = 5e - 2$:
\[
    1.05 y_8 - 2 y_9 = 0.03 - 0.95(5e - 2)
\]

\textbf{Алгоритм метода прогонки:}

\textbf{Прямой ход (прогонка):}
\begin{align*}
    \alpha_1 & = -\frac{c_1}{b_1}                                                             \\
    \beta_1  & = \frac{d_1}{b_1}                                                              \\
    \alpha_i & = -\frac{c_i}{b_i + a_i \alpha_{i-1}}, \quad i = 2, \ldots, m                  \\
    \beta_i  & = \frac{d_i - a_i \beta_{i-1}}{b_i + a_i \alpha_{i-1}}, \quad i = 2, \ldots, m
\end{align*}

\textbf{Обратный ход (вычисление решения):}
\begin{align*}
    y_m & = \beta_m                                              \\
    y_i & = \alpha_i y_{i+1} + \beta_i, \quad i = m-1, \ldots, 1
\end{align*}

После получения $y_1, \ldots, y_9$, добавляем граничные значения $y_0 = 6$ и $y_{10} = 5e - 2$.

Результаты вычисления представим в виде таблицы значений $x_i$ и приближённого решения $\tilde{y}_i$ с последующим сравнением с точным аналитическим решением $y(x) = 1 + 5e^x - 3x$.

%%%%%%%%%%%%%%%%%%%%%%%
%%% Метод стрельбы %%%%
%%%%%%%%%%%%%%%%%%%%%%%
\subsection*{Метод стрельбы}
\addcontentsline{toc}{subsection}{Метод стрельбы}

Рассмотрим краевую задачу
\[
    y'' + p(x)y' + q(x)y = f(x), \qquad y(a)=A, \qquad y(b)=B.
\]

В нашем случае:
\[
    y'' - y' = 3, \qquad y(0)=6, \qquad y(1)=5e-2.
\]
Следовательно,
\[
    p(x)=-1, \qquad q(x)=0, \qquad f(x)=3, \qquad a=0, \qquad b=1,
\]
\[
    A=6, \qquad B=5e-2.
\]

Согласно методу стрельбы, решение краевой задачи представляется в виде
\[
    y(x)=y_0(x)+C_1y_1(x),
\]
где \(y_0(x)\) $-$ решение неоднородного уравнения, а \(y_1(x)\) $-$ решение соответствующего однородного уравнения.

В сеточном варианте метода отрезок \([a,b]\) разбивается на \(n\) частей:
\[
    x_i=a+ih, \qquad i=0,1,\dots,n, \qquad h=\frac{b-a}{n}.
\]
При \(n=10\) получаем
\[
    h=\frac{1-0}{10}=0.1, \qquad x_i=0.1i, \qquad i=0,1,\dots,10.
\]

Производные заменяются центральными разностями:
\[
    y_i' \approx \frac{y_{i+1}-y_{i-1}}{2h},
    \qquad
    y_i'' \approx \frac{y_{i+1}-2y_i+y_{i-1}}{h^2}.
\]

Подставляя эти выражения в уравнение
\[
    y'' + p_i y' + q_i y = f_i,
\]
получаем разностную схему:
\[
    \frac{y_{i+1}-2y_i+y_{i-1}}{h^2}
    + p_i \frac{y_{i+1}-y_{i-1}}{2h}
    + q_i y_i
    = f_i,
    \qquad i=1,2,\dots,n-1.
\]

Умножая на \(h^2\), имеем
\[
    y_{i+1}-2y_i+y_{i-1}
    +\frac{p_i h}{2}(y_{i+1}-y_{i-1})
    +q_i h^2 y_i
    = f_i h^2.
\]

После группировки членов получаем рекуррентную формулу:
\[
    \left(1-\frac{p_i h}{2}\right)y_{i-1}
    +\left(q_i h^2-2\right)y_i
    +\left(1+\frac{p_i h}{2}\right)y_{i+1}
    = f_i h^2,
\]
откуда
\[
    y_{i+1}
    =
    \frac{
    f_i h^2 + (2-q_i h^2)y_i - \left(1-\frac{p_i h}{2}\right)y_{i-1}
    }{
    1+\frac{p_i h}{2}
    }.
\]

\textbf{Построение двух вспомогательных решений.}

\begin{enumerate}
    \item \(y_0[i]\) $-$ решение неоднородного уравнения
          \[
              y'' + p(x)y' + q(x)y = f(x)
          \]
          с начальными сеточными условиями
          \[
              y_0[0]=A, \qquad y_0[1]=D_0.
          \]

    \item \(y_1[i]\) $-$ решение однородного уравнения
          \[
              y'' + p(x)y' + q(x)y = 0
          \]
          с начальными сеточными условиями
          \[
              y_1[0]=0, \qquad y_1[1]=D_1, \qquad D_1 \neq 0.
          \]
\end{enumerate}

Тогда искомое решение ищется в виде
\[
    y[i]=y_0[i]+C_1 y_1[i], \qquad i=0,1,\dots,n.
\]

Константа \(C_1\) находится из правого граничного условия:
\[
    y[n]=B.
\]
Следовательно,
\[
    y_0[n]+C_1 y_1[n]=B,
\]
откуда
\[
    C_1=\frac{B-y_0[n]}{y_1[n]}.
\]

\textbf{Итоговые формулы}

Так как
\[
    p_i=-1, \qquad q_i=0, \qquad f_i=3,
\]
то рекуррентные формулы принимают вид:

для неоднородного уравнения
\[
    y_0[i+1]
    =
    \frac{
    3h^2 + 2y_0[i] - \left(1+\frac{h}{2}\right)y_0[i-1]
    }{
    1-\frac{h}{2}
    },
    \qquad i=1,2,\dots,n-1,
\]

для однородного уравнения
\[
    y_1[i+1]
    =
    \frac{
    2y_1[i] - \left(1+\frac{h}{2}\right)y_1[i-1]
    }{
    1-\frac{h}{2}
    },
    \qquad i=1,2,\dots,n-1.
\]

При \(h=0.1\) имеем
\[
    1-\frac{h}{2}=0.95, \qquad 1+\frac{h}{2}=1.05, \qquad 3h^2=0.03.
\]

Поэтому получаем:

для неоднородного уравнения
\[
    y_0[i+1]=\frac{0.03+2y_0[i]-1.05\,y_0[i-1]}{0.95},
    \qquad i=1,\dots,9,
\]

для однородного уравнения
\[
    y_1[i+1]=\frac{2y_1[i]-1.05\,y_1[i-1]}{0.95},
    \qquad i=1,\dots,9.
\]

\textbf{Выбор начальных сеточных значений.}

Согласно используемому алгоритму, берём
\[
    D_0=A,
    \qquad
    D_1=h.
\]
Тогда
\[
    y_0[0]=6, \qquad y_0[1]=6,
\]
\[
    y_1[0]=0, \qquad y_1[1]=0.1.
\]

После последовательного вычисления значений \(y_0[i]\) и \(y_1[i]\) для \(i=2,\dots,10\) находим коэффициент
\[
    C_1=\frac{B-y_0[10]}{y_1[10]},
    \qquad B=5e-2.
\]
